\documentclass{article}

\usepackage{amsmath}

\title{Ch. 30 CW Solutions}
\author{Principles of Macroeconomics}
\date{\today}

\begin{document}

\maketitle

\begin{enumerate}

\item

	\begin{enumerate}

	\item Reserves decrease by \$1 million. That means deposits decrease by
		\begin{equation*}
		\$1 \text{ million} \cdot \frac{1}{0.1} = \$10 \text{ million}.
		\end{equation*}
		There's no change in currency, so, the money supply decreases by \$10 million.
		
	\item Banks might choose to hold excess reserves if the Fed pays a high enough interest rate on reserves. They also might want to hold excess reserves if they anticipate heavy withdrawals. 
	
	The money multiplier stays the same because the overall reserve ratio remains 10\%. Since the money multiplier remains the same, the money supply doesn't change. 
		
	\end{enumerate}
	
\item 

	\begin{enumerate}
	
	\item The money multiplier is \[ MM = \frac{1}{0.1} = 10. \] The money supply is
		\begin{equation*} 
		S_M = C + D = 0 + res \cdot MM = 0 + \$100 \text{ billion} \cdot 10 = \$1 \text{ trillion}
		\end{equation*}
		
	\item Since households hold no currency, banks can't increase their reserves. As a result, reserves stay constant at \$100 billion. The new money multipliers is $1 / 0.2 = 5$, so the money supply becomes 
		\begin{equation*}
		\$100 \text{ billion} \cdot 5 = \$500 \text{ billion}. 
		\end{equation*}
	
	\end{enumerate}
	
\item

	\begin{enumerate}
	
	\item It will buy bonds because buying bonds requires the Fed to pay bond owners. Paying bond owners allows the Fed to introduce new dollars into circulation.
	
	\item Let $x$ be reserves. The Fed will need to increase reserves by
		\begin{gather*}
		x \cdot \frac{1}{0.2} = \$40 \text{ million} \\
		5x = \$40 \text{ million} \\
		x = \$8 \text{ million}
		\end{gather*}
		That means the Fed will need to introduce \$8 million dollars worth of new dollars to the economy. It should buy \$8 million worth of bonds. 
	
	\end{enumerate}
	
\item 

	\begin{enumerate}
	
	\item $S_M = C + D = \$2,000 + \$0 = \$2,000$
	
	\item $S_M = C + D = \$0 + res \cdot MM = 0 + \$2,000 \cdot \frac{1}{1} = \$2,000$
	
	\item We get the following system:
	
		\begin{equation*}
		\begin{cases}
		C = D \\
		D = res \cdot MM = (\$2,000 - C) \cdot \frac{1}{1} = \$2,000 - C
		\end{cases}
		\end{equation*}
		
		By substitution,
		
		\begin{gather*}
		C = \$2,000 - C \\
		2C = \$2,000 \\
		C = \$1,000 \\
		D = \$1,000 \\
		\end{gather*}
		
		Therefore, \[S_M = C + D = \$1,000 + \$1,000 = \$2,000.\]
		
	\item $S_M = C + D = 0 + res \cdot MM = 0 + \$2,000 \cdot \frac{1}{0.1} = \$20,000$
	
	\item We get the following system:
		\begin{equation*}
		\begin{cases}
		C = D \\
		D = res \cdot MM = (\$2,000 - C) \cdot \frac{1}{0.1} = \$20,000 - 10C
		\end{cases}
		\end{equation*}
		Then by substitution, we get
		\begin{gather*}
		C = \$20,000 - 10C \\
		11C = \$20,000 \\
		C = \$1,818.18 \\
		D = \$1,818.18 \\
		\end{gather*}
		
		Therefore, \[S_M = C + D = \$1,818.18 + \$1,818.18 = \$3,636.36.\]
	
	\end{enumerate}
	
\item

	\begin{enumerate}

	\item 
		\begin{gather*}
		D = res \cdot MM = \$200 \cdot \frac{1}{0.25} = \$800 \\
		S_M = C + D = \$200 + \$800 = \$1,000
		\end{gather*}
		
	\item We get the following system:
		\begin{equation*}
		\begin{cases}
		C = res \\
		\$650 = C + res \cdot MM = C + res \cdot \frac{1}{0.25} = C + 4 \cdot res
		\end{cases}
		\end{equation*}
		Then by substitution, we get
		\begin{gather*}
		\$650 = res + 4 \cdot res \\
		\$650 = 5 \cdot res \\
		\$130 = res \\
		C = \$130
		\end{gather*}
		
	\item First, we can find deposits:
		\begin{gather*}
		S_M = C + D \\
		\$1,500 = \$400 + D \\
		\$1,100 = D
		\end{gather*}
		Then, we can find the reserve ratio:
		\begin{equation*}
		\text{reserve ratio} = \frac{\$110}{\$1,100} = 0.10
		\end{equation*}
		

	\end{enumerate}

\end{enumerate}

\end{document}